% Figure: van der Waals isotherms below the critical temperature, showing
% the non-monotonic loop and the Maxwell equal-area construction that
% replaces the loop with a horizontal coexistence line. Plotted from the
% reduced van der Waals equation p_r = 8 T_r/(3 v_r - 1) - 3/v_r^2, which
% is well-defined for v_r > 1/3. We clamp the domain to v_r >= 0.5 to keep
% the divergence near v_r=1/3 off the plotted range.
\begin{figure}[htbp]
\centering
\begin{tikzpicture}
\begin{axis}[
  width=12cm, height=8cm,
  xlabel={reduced volume $v_r=v/v_c$},
  ylabel={reduced pressure $p_r=p/p_c$},
  xmin=0.4, xmax=4.0, ymin=0, ymax=1.6,
  axis lines=left,
  tick label style={font=\scriptsize},
  label style={font=\small},
  legend pos=north east,
  legend style={font=\scriptsize},
  domain=0.5:4.0, samples=200,
  clip=true,
]
% subcritical reduced isotherm T_r = 0.90, showing the van der Waals loop
\addplot[engred, very thick] {8*0.90/(3*x - 1) - 3/(x*x)};
\addlegendentry{$T_r=0.90$ (vdW loop)}
% critical isotherm T_r = 1.0
\addplot[engorange, thick] {8*1.0/(3*x - 1) - 3/(x*x)};
\addlegendentry{$T_r=1.00$}
% supercritical isotherm T_r = 1.10
\addplot[enggray, thick] {8*1.10/(3*x - 1) - 3/(x*x)};
\addlegendentry{$T_r=1.10$}
% Maxwell horizontal coexistence line for T_r=0.90.
% The equal-area saturation pressure is p_r approx 0.647 at T_r=0.90.
\addplot[engblue, very thick, dashed] coordinates {(0.60,0.647) (3.5,0.647)};
\addlegendentry{Maxwell tie line}
% liquid and vapour roots on the tie line (approximate)
\addplot[only marks, mark=*, engblue, mark size=1.8pt] coordinates {(0.60,0.647) (3.45,0.647)};
\node[engblue, font=\scriptsize, anchor=south east] at (axis cs:0.60,0.68) {$v_f$};
\node[engblue, font=\scriptsize, anchor=south west] at (axis cs:3.45,0.68) {$v_g$};
\end{axis}
\end{tikzpicture}
\caption[Van der Waals isotherms and the Maxwell construction]{Reduced
van der Waals isotherms $p_r = 8T_r/(3v_r-1) - 3/v_r^2$. Above the
critical temperature ($T_r=1.10$, grey) the isotherm is monotonic. At
the critical temperature ($T_r=1.00$, orange) it has a horizontal
inflexion. Below the critical temperature ($T_r=0.90$, red) it develops a
non-monotonic loop with a region of $(\partial p/\partial v)_T>0$, which
is mechanically unstable and therefore unphysical. The Maxwell
equal-area construction replaces the loop with the horizontal coexistence
line (blue dashed) drawn so that the two lobes the line cuts from the
loop have equal area; the endpoints give the saturated-liquid volume
$v_f$ and the saturated-vapour volume $v_g$ at that temperature. The
equation is plotted only for $v_r\ge 0.5$ to keep the divergence at
$v_r=1/3$ off the range.}
\label{fig:vol04ch02:vdw-loops}
\end{figure}
